\section{結果および考察}

\subsection{アンケート結果および考察}

\subsubsection{生成されたプロットに対する満足度}
\begin{figure}[htbp]
\begin{center}
\includegraphics[width=70mm]{img/5/quality2.png}
\caption{本アプリケーションが生成するプロットの各項目における満足度を表すアンケート結果.}
\label{quality}
\end{center}
\end{figure}
図\ref{quality}より, 本アプリケーションは, 目新しさの提供に課題がある一方で, 幅広さや面白さといった観点で, プロのストーリークリエーターから見ても創作活動に役立つものであったことがわかる.

複合的な作品である『ブラック・ジャック』の分析データを, プロットの生成に活用したことが, LLM に幅広い言葉の出力を促したのではないかと考えられる.

一方, 目新しさの課題については, 出力が学習データに制限される LLM を創造的活動に用いる上で重要な課題の一つであるといえる.

\subsubsection{LLM との共創における貢献度}
\begin{figure}[htbp]
\begin{center}
\includegraphics[width=70mm]{img/5/creativity_control.png}
\caption{本アプリケーションにおける LLM との共創を促す貢献について評価されたアンケート結果.}
\label{qc}
\end{center}
\end{figure}
図\ref{qc}より, 作業時間の短縮に関する貢献を除き, 低い評価を受けていないことがわかる. 特に, 使いやすさと指示の反映に関する項目について高い評価を受けており, 本アプリケーションが LLM との創作活動を向上させる可能性を示唆した結果であるといえる.

作業時間については, LLM からのレスポンスの遅さに不満の声が聞かれ, より高速に応答する LLM がこの問題を解決すると考えられる.

創造性への刺激については, まずまずの評価を得ているが, 目新しさを提供する課題の取り組みにより, さらなる貢献をもたらすことができると考える.

\subsection{ログの結果および考察}

\subsubsection{作業時間}
プロット制作という創造的活動の作業時間を正確に把握することは難しい. 実際に手を動かしていなくとも, アイデアや構想を練ることなどに時間は使われるからである. 一方で創作活動の時間効率性を高めることは, 創作支援の重要な項目の一つであると考えられる. よって本論文では, 本アプリケーションの操作ログより作業時間を概算し, 本アプリケーションの時間効率への寄与について考察する.
\begin{figure}[htbp]
\includegraphics[width=60mm]{img/5/time_bar_A1A2.png}
\caption{ボタンの押下間隔から算出した作業時間. 被験者グループAは制作期間内において2つのプロットを制作したため作品ごとに作業時間を集計した. Average1 は全体平均, Average2 は外れ値を除いた平均値を表す.}
\label{time}
\end{figure}

作業時間は, ログデータとして集計した, 本アプリケーション画面のボタン押下間隔のうち, 98 パーセンタイルから 99 パーセンタイルの間に位置する 1 時間を区切りとして算出した. つまり, ボタン押下間隔が 1 時間以上開いている値は作業時間に含めずに, ボタン押下間隔が 1 時間未満の値を全て足し上げ算出した値を便宜的に作業時間とみなすこととした.

図\ref{qc}で示した, 「作業時間が短くなったか」 というアンケートへの回答のうち「短くなったとは思わない」と回答したのは, グループ C と E であり, そのほかのグループは「短くなったと思う」と回答している. 図\ref{time}からわかる通り, グループ C は作業時間を最も費やしたグループである一方, グループ E は最も作業時間の少ないグループであった.

グループ E が作業時間の短縮を感じなかったのは, LLM の応答の遅さや, 自身のプロット制作にかかる時間が元々十分に早いことなどが理由であると考えられる. 

\begin{table}[hbtp]
\small
\begin{center}
{\caption{各グループの編集機能の利用数}
\label{edit_action}
\begin{tabular}{ccccc} \hline
A & B & C & D & E  \\ \hline
116 & 60 & 176 & 33 & 45\\
\hline
\end{tabular}
}
\end{center}
\end{table}
\vspace{-1.5zh}

表\ref{edit_action}に示した通り, グループ C は最も多く編集機能を用いたグループでもある. これより, 編集機能における操作性などが時間効率に影響していると考えられる.

一方で, 作業時間の短縮を感じたと回答したグループはいずれも作業時間が 10 時間未満であり, このことに本アプリケーションの時間効率への寄与があったと考えられる.

\subsubsection{編集機能により生成したプロットの採用率}


\begin{table}[hbtp]
\small
\begin{center}
{\caption{グループごとの編集機能により生成したプロットの採用率とそれらの平均}
\label{adopt_t}
\begin{tabular}{cccccc} \hline
A & B & C & D & E & mean \\ \hline
0.71 & 0.91 & 0.70 & 0.50 & 0.63 & 0.69\\
\hline
\end{tabular}
}
\end{center}
\end{table}
表\ref{adopt_t}より, 編集機能を用いて, 既存のプロットを新たなプロットへと書き換えた際, ほとんどのグループで新旧プロットのうち新しいプロットを採用する割合が大きいことがわかる. 各グループの採用率の平均もおよそ 70 \% と高い割合を示している. 

この結果は, 表\ref{qc}で示したアンケート結果, 「編集機能により面白さは増加したか」という問において高い評価が得られた結果と符合している. つまり, これらの結果は, LLM とのインタラクティブなプロット創作を可能とする本アプリケーションの編集機能が, プロットの質向上に寄与することを示唆する結果であるといえる.

% \begin{figure}[htbp]
% \begin{center}
% \includegraphics[width=70mm]{img/5/comparison.png}
% \caption{編集機能により生成したプロットの採用・不採用数}
% \label{adopt}
% \end{center}
% \end{figure}

\subsubsection{用いられた編集機能の割合}

\begin{figure}[htbp]
\includegraphics[width=65mm]{img/5/edit_action_pie_chart_N.png}
\caption{グループごとに正規化した各編集機能の利用率}
\label{edit_action}
\end{figure}

図\ref{edit_action}より, 編集機能の中でフリーワードでやり取りする機能が最も多く利用され, グループごとに正規化した利用率の中でも50\%以上を占めることがわかる. 

この結果から, 生成されたプロットが, 型にはまらない自由な発想を引き出した可能性や, プロのストーリークリエーターが自身の言葉による制御を望んていたこと, プロットの多様性により, 機能を限定しない編集が必要となっていたことなどが考えられる.

フリーワードでのやり取りに対する支援の拡充が求められていると考える.